\section{Introdução}

Este relatório descreve o desenvolvimento de \textbf{SportFlow}, uma aplicação web de comércio eletrónico para artigos desportivos, desenvolvida no âmbito da unidade curricular de Sistemas Distribuídos da Universidade da Beira Interior.

\subsection{Motivação}

O comércio eletrónico tem registado um crescimento expressivo nos últimos anos, com particular destaque para o segmento de artigos desportivos. A SportFlow surge como resposta à necessidade de uma plataforma completa que permita a um utilizador descobrir, pesquisar e adquirir equipamento desportivo de forma intuitiva, e a um administrador gerir o catálogo, stock e vendas.

\subsection{Objetivos}

Os principais objetivos do projeto são:

\begin{itemize}
    \item Desenvolver uma API REST robusta com autenticação por token JWT;
    \item Implementar um frontend moderno e responsivo em React/TypeScript;
    \item Garantir a persistência de dados em MySQL com migrações versionadas (Flyway);
    \item Aplicar princípios de segurança (BCrypt, HTTPS-ready, CORS configurado);
    \item Separar claramente as responsabilidades entre cliente, servidor e base de dados, seguindo a arquitetura de três camadas típica de sistemas distribuídos.
\end{itemize}

\subsection{Estrutura do Relatório}

O relatório está organizado da seguinte forma: a Secção~\ref{sec:design} apresenta o modelo de dados e a arquitetura do sistema; a Secção~\ref{sec:impl} detalha a implementação das funcionalidades principais; a Secção~\ref{sec:demo} demonstra o funcionamento da aplicação; e a Secção~\ref{sec:conc} apresenta as conclusões e trabalho futuro.
